\section{BMHV projectors and dimension recurrences}
\label{app:bmhv-dimension-shift}

The BMHV decomposition is orthogonal,
\begin{equation}
 \bar g^{\mu\rho}\widehat g_{\rho}{}^{\nu}=0,
 \qquad
 \bar g^\mu{}_{\mu}=4,
 \qquad
 \widehat g^\mu{}_{\mu}=-2\eps.
 \label{eq:bmhv-orthogonal-metrics}
\end{equation}
For an integration momentum $\ell$, define
\(\mu_\ell^2=-\widehat\ell^2\).  With the normalized measure in
Eq.~\eqref{eq:ibp-measure}, a one-loop scalar integrand whose external vectors
are four dimensional obeys
\begin{equation}
 \int\frac{\dd^D\ell}{i\pi^{D/2}}
 (\mu_\ell^2)^r F(\ell)
 =(-\eps)_r
 \int\frac{\dd^{D+2r}\ell}{i\pi^{(D+2r)/2}}F(\ell),
 \qquad r\in\mathbb Z_{\geq0},
 \label{eq:one-loop-dimension-shift}
\end{equation}
where $(a)_r=\Gamma(a+r)/\Gamma(a)$.  The shifted integral contains the same
denominator polynomials and cut slots, interpreted in $D+2r$ dimensions.

At several loops, mixed contractions
${\widehat\ell_i\cdot\widehat\ell_j}$ are generated as well as powers of
$\mu_{\ell_i}^2$.  Introduce Schwinger parameters for the complete family.
The Gaussian integral over the evanescent components expresses every rank
(2r) monomial as a polynomial in cofactors of the loop quadratic-form matrix
and gives scalar integrals in (D+2r) dimensions.  Symbolically,
\begin{equation}
 I^{(D)}_{\mathcal F}[P_{2r}(\widehat\ell)]
 =\mathcal D_{P_{2r}}(\boldsymbol\nu,D),
 I_{\mathcal F}^{(D+2r)}(\boldsymbol\nu),
 \label{eq:multiloop-evanescent-shift}
\end{equation}
where $\mathcal D_{P_{2r}}$ is the exact Tarasov raising operator for the
chosen topology.  It changes propagator powers according to the same Schwinger
polynomial; it is not a universal replacement of
$\widehat\ell_i\cdot\widehat\ell_j$ by a number.

Tarasov recurrences then lower the dimension,
\begin{equation}
 I_{\mathcal F}^{(D+2r)}(\boldsymbol\nu)
 =\sum_{a}T_a(\boldsymbol\nu,D,\{s_{ij}\})
 I_{\mathcal F}^{(D)}(\boldsymbol\nu+\Delta\boldsymbol\nu_a),
 \label{eq:tarasov-lowering}
\end{equation}
with rational $T_a$.  Terms in which a mandatory cut obtains nonpositive
power vanish by Eq.~\eqref{eq:cut-family-condition}.  The recurrence is
checked by substituting the shifted-dimension representation back into the
Schwinger-parametric identity and by verifying that no evanescent scalar
product remains.  Coefficients originating in the original Dirac trace stay
at $D=4-2\eps$; only the scalar integral in
Eq.~\eqref{eq:multiloop-evanescent-shift} is dimension shifted.
